\section{Discussion}\label{sec:disc}

Before implementing prediction, one can first ask whether its lead changes what recovery can accomplish. Prediction is most plausible when recovery is slow enough for an earlier trigger to cross the boundary in Proposition~\ref{prop:boundary}, and when degradation rates vary in a way that can be estimated before that boundary. At one known constant rate, a fixed pad reproduces the studied predictive trigger. Under fast, committed serialization, tuning may remove predictive lead altogether.

The simulators produce different results. In the single-margin study, the exploratory figures favor prediction over the unpadded rules and, in a few settings, over the in-sample fixed pad. After equal-budget, seed-disjoint tuning, the predictive arm intervenes less often on clear paths and has a lower descriptive successful-only elapsed time. It also records two held-out conflict collisions where the fixed arm records none, so an efficiency comparison at matched outcomes is unavailable. In the coordination study, both tuned arms succeed throughout the held-out conflict grid and the predictive arm finishes earlier. Its lower rate of clear-path serialization is consistent with adaptation to degradation rate.

The coordination result depends on the recovery design. The first trigger commits one total order, after which blockers and release events determine the motion. Prediction changes when this commitment occurs; the order is never revised. The paired interval thus describes variation over measurement noise on the fixed 27-scenario grid. It says nothing about unseen geometries, robot counts, dynamics, or online re-planning.

Proposition~\ref{prop:reduction} also suggests an intermediate design: scale a current-margin threshold using an online rate estimate, without constructing a larger predict-then-recover stack. Its practical value would still depend on estimator delay, noise, and changes in degradation rate.

Zero observed collisions is not a safety guarantee. The pointwise Wilson intervals are not simultaneous bounds over the scenario family, and holding out noise seeds does not create a new population of robot environments. The v1 comparison accordingly tunes the simpler early-trigger baseline, verifies that untreated paths conflict, separates tuning from testing, and reports failures before conditional timing.
